\subsection{FLAGS and STATUS Registers}\label{section:state-register-formats}
FLAGS and STATUS are accessed through dedicated read/write instructions.

\begin{manuallistedformatdiagram}{FLAGS Register Format}
\manualformatrow{FLAGS[15:0]}{%
\manualformatreserved{reserved}{12}
\manualformatfield{Z}{1}
\manualformatfield{N}{1}
\manualformatfield{C}{1}
\manualformatfield{V}{1}
}
\end{manuallistedformatdiagram}

\manualtablecaption{FLAGS Fields}
\begin{manualtabular}{@{}>{\raggedright\arraybackslash}p{0.75in}>{\raggedright\arraybackslash}p{0.65in}>{\raggedright\arraybackslash}p{4.20in}@{}}
\toprule
\rowcolor{ManualHeaderFill}
\textbf{Field} & \textbf{Bits} & \textbf{Meaning}\\
\midrule
reserved & \texttt{15..4} & reads as zero and must be zero when written\\
\texttt{Z} & \texttt{3} & result value is zero\\
\texttt{N} & \texttt{2} & most significant result bit at the operand size\\
\texttt{C} & \texttt{1} & unsigned carry out for addition or unsigned borrow for subtraction\\
\texttt{V} & \texttt{0} & signed overflow or an explicitly reported exceptional condition\\
\bottomrule
\end{manualtabular}

\hyperref[instr:rdflags]{\texttt{RDFLAGS}} and \hyperref[instr:wrflags]{\texttt{WRFLAGS}} are unprivileged.
\texttt{RDFLAGS} reads and zero-extends the 16-bit image; \texttt{WRFLAGS} writes the low 16 bits and requires every
reserved bit to be zero. Instruction-specific flag production and conditional use are defined in
\hyperref[page:flags-and-condition-codes]{Flags and Condition Codes}.

\begin{manuallistedformatdiagram}{STATUS Register Format}
\manualformatrow{STATUS[15:0]}{%
\manualformatreserved{reserved}{10}
\manualformatfield{I}{1}
\manualformatfield{P}{1}
\manualformatfield{R}{1}
\manualformatfield{T}{1}
\manualformatfield{N}{1}
\manualformatfield{E}{1}
}
\end{manuallistedformatdiagram}
{\small
In the diagram, \texttt{I}, \texttt{P}, \texttt{R}, \texttt{T}, \texttt{N}, and \texttt{E}
abbreviate \texttt{IE}, \texttt{PM}, \texttt{RF}, \texttt{TF}, \texttt{NI}, and \texttt{EA}.
\par}

\manualtablecaption{STATUS Fields}
\begin{manualtabular}{@{}>{\raggedright\arraybackslash}p{0.75in}>{\raggedright\arraybackslash}p{0.65in}>{\raggedright\arraybackslash}p{4.20in}@{}}
\toprule
\rowcolor{ManualHeaderFill}
\textbf{Field} & \textbf{Bits} & \textbf{Meaning}\\
\midrule
reserved & \texttt{15..6} & reads as zero and must be zero when written\\
\texttt{IE} & \texttt{5} & permits maskable-interrupt admission when the other admission conditions hold\\
\texttt{PM} & \texttt{4} & selects user mode when zero and supervisor mode when one\\
\texttt{RF} & \texttt{3} & one-shot suppression of \texttt{DEBUG\_TRACE}\\
\texttt{TF} & \texttt{2} & requests post-success \texttt{DEBUG\_TRACE} for a trace unit when \texttt{RF} is clear\\
\texttt{NI} & \texttt{1} & inhibits NMI admission while preserving the NMI pending\\
\texttt{EA} & \texttt{0} & hardware-managed architectural event-active state\\
\bottomrule
\end{manualtabular}

\hyperref[instr:rdstatus]{\texttt{RDSTATUS}} is unprivileged; \hyperref[instr:wrstatus]{\texttt{WRSTATUS}} is
supervisor-only. \texttt{WRSTATUS} atomically updates \texttt{IE}, \texttt{NI}, \texttt{TF}, and \texttt{RF};
reserved bits must be zero, and supplied \texttt{PM} and \texttt{EA} must equal their current values. Privilege
transitions are defined in the \hyperref[page:privileged-execution-model]{Privileged Execution Model}; event
admission, tracing, and event-active transitions are defined in the
\hyperref[page:architectural-event-processing-model]{Architectural Event Processing Model}.

\subsection{Floating-Point Register Model}\label{section:floating-point-register-model}
The floating-point extension exposes F0 through F15 as 64-bit registers when the floating-point
instruction group is implemented.

\subsection{FFLAGS and FSTATUS Registers}
FFLAGS holds accrued IEEE-754 exception flags. FSTATUS holds the matching trap enables, rounding mode, and optional
non-default floating-point modes. Both registers are present only with the floating-point extension and reset to zero.

\begin{manuallistedformatdiagram}{FFLAGS Register Format}
\manualformatrow{FFLAGS[15:0]}{%
\manualformatreserved{reserved}{11}
\manualformatfield{I}{1}
\manualformatfield{Z}{1}
\manualformatfield{O}{1}
\manualformatfield{U}{1}
\manualformatfield{X}{1}
}
\end{manuallistedformatdiagram}

\begin{manuallistedformatdiagram}{FSTATUS Register Format}
\manualformatrow{FSTATUS[15:0]}{%
\manualformatreserved{reserved}{6}
\manualformatfield{D}{1}
\manualformatfield{A}{1}
\manualformatfield{F}{1}
\manualformatfield{R}{2}
\manualformatfield{I}{1}
\manualformatfield{Z}{1}
\manualformatfield{O}{1}
\manualformatfield{U}{1}
\manualformatfield{X}{1}
}
\end{manuallistedformatdiagram}
{\small
In the diagrams, \texttt{I}, \texttt{Z}, \texttt{O}, \texttt{U}, and \texttt{X} denote the
\texttt{NV}, \texttt{DZ}, \texttt{OF}, \texttt{UF}, and \texttt{NX} flag or enable bits.
In FSTATUS, \texttt{D}, \texttt{A}, \texttt{F}, and \texttt{R} denote
\texttt{DN}, \texttt{DAZ}, \texttt{FTZ}, and \texttt{RM}.
\par}

\manualtablecaption{Floating-Point State Fields}
\begin{manualtabular}{@{}>{\raggedright\arraybackslash}p{0.75in}>{\raggedright\arraybackslash}p{0.65in}>{\raggedright\arraybackslash}p{3.95in}@{}}
\toprule
\rowcolor{ManualHeaderFill}
\textbf{Field} & \textbf{Bits} & \textbf{Meaning}\\
\midrule
\texttt{NV..NX} & \texttt{4..0} & accrued invalid, divide-by-zero, overflow, underflow, and inexact flags\\
\texttt{NV\_EN..NX\_EN} & \texttt{4..0} & trap enables corresponding to FFLAGS.NV through FFLAGS.NX\\
\texttt{RM} & \texttt{6..5} & 0 nearest-even; 1 toward zero; 2 toward negative; 3 toward positive\\
\texttt{FTZ} & \texttt{7} & flush tiny results to signed zero\\
\texttt{DAZ} & \texttt{8} & treat subnormal inputs as signed zero\\
\texttt{DN} & \texttt{9} & produce the architectural default NaN\\
reserved & \texttt{15..10} & must be zero\\
\bottomrule
\end{manualtabular}
